\chapter{Introduction} % Main chapter title
\label{Chapter1}

Biodiversity loss is an escalating global concern, with the Arctic tundra particularly vulnerable to environmental changes. They are driving rapid ecological shifts in this sensitive biome, altering plant community structures, species composition, and habitat distribution (\cite{garcia_criado_plant_2025}; \cite{bjorkman_status_2019}; \cite{wenzl_vegetation_2024}). Monitoring these transformations is essential for understanding and responding to ongoing ecosystem changes. However, the Arctic’s vast, remote, and often inhospitable terrain makes traditional in-situ biodiversity monitoring highly impractical.

Remote sensing offers a powerful alternative for large-scale ecological assessment. Spectral diversity, variation in reflected light across wavelengths, has shown strong potential as a proxy for biodiversity, especially in regions with high species turnover (\cite{elisa_van_cleemput_making_2022}; \cite{pijnenburg_mapping_2023}). Yet, questions remain about how effectively spectral diversity captures ecological variation in Arctic ecosystems, where vegetation is low-growing and plant canopies are small. Climate change further complicates this, as it drives divergent responses among plant groups: vascular plants (e.g., deciduous shrubs and evergreens) often expand with warming, while non-vascular species (e.g., lichens and bryophytes) tend to decline (\cite{garcia_criado_plant_2025}; \cite{bjorkman_status_2019}). These contrasting dynamics underscore the need for scalable, sensitive tools capable of detecting multidimensional biodiversity shifts.

The \ac{svh} posits that spectral heterogeneity in remote sensing data correlates positively with ecological diversity (\cite{palmer_quantitative_2002}; \cite{schmidtlein_spectral_2017}). However, empirical support for the \ac{svh} has been mixed. As \cite{palmer_quantitative_2002} noted, the relationship between habitat heterogeneity and species richness is highly context dependent, and areas with high species richness may not always exhibit high spectral variability. This inconsistency is especially pronounced at coarse spatial scales and is further influenced by differences in study design, including spatial extent, image resolution, spectral granularity, and timing of data acquisition (\cite{rocchini_remotely_2010}).

Due to the Arctic’s harsh conditions and logistical challenges, ground-based validation data are limited. This has motivated the development of tools like the R package biodivMapR \cite{feret_biodivmapr_2020}, which enables $\alpha$- and $\beta$-diversity mapping from hyperspectral imagery using the concept of spectral species—clusters of pixels with similar spectral reflectance. A critical parameter in this process is the number of spectral species, which many studies have set arbitrarily (\cite{feret_mapping_2014}; \cite{pijnenburg_mapping_2023}; \cite{zhang_assessment_2023}; \cite{robertson_effects_2023}; \cite{liccari_use_2022}). In this study, I applied a data-driven approach using \ac{wcss} metrics to determine the optimal number of spectral species for each hyperspectral image, enhancing the ecological relevance of the clustering. I investigated whether spectral diversity derived from airborne hyperspectral imagery could serve as a proxy for three key metrics of Arctic tundra biodiversity: habitat type diversity, plant community diversity, and taxonomic diversity. Specifically, I used high-resolution imagery from \ac{nasa}’s \ac{avirisng} sensor (425 spectral bands at ~5 m resolution) in combination with vegetation records from the \ac{ava}. Spectral species were defined using biodivMapR with k-means clustering, parameterized through \ac{wcss}-based optimization. To my knowledge, this is the first study to apply such a data-driven parameterization within the biodivMapR framework.

Given the coarse spatial resolution of the imagery relative to the small stature of Arctic vegetation, I hypothesized that only broad-scale diversity patterns, such as those at the habitat or plant community level, would be reliably captured. Through this approach, I aimed to evaluate spectral diversity as a practical and robust tool for biodiversity monitoring in Arctic ecosystems. The findings contribute to advancing remote sensing applications for long-term ecological assessment in the context of rapid environmental change.